\section{\texorpdfstring{Coupling independence for \(q\ge\Deg+3\) at large girth}
{Coupling independence for q >= Delta + 3 at large girth}}
\label{sec:high-girth-soft}

\begin{theorem}[Large-girth coupling independence]
\label{thm:high-girth-soft-ci}
Let $q,\Deg$ be integers with $\Deg\ge3$ and $q\ge\Deg+3$.  There exist
$g_*=g_*(q,\Deg)\ge3$ and $C_*=C_*(q,\Deg)<\infty$ such that, for every
$G\in\Gdeg$ and every pinning $\tau$ with
\(\operatorname{girth}(G^\tau)\ge g_*\), the bound
below holds for every $x\in[0,1]$, every $r\in V^\tau$, and all
$a,b\in\colours$:
\[
  \WHam\!\left(
    \mupin{G}{\tau^{r=a}}{x},
    \mupin{G}{\tau^{r=b}}{x}
  \right)\le C_*.
\]
The two laws live on the common free set $V^\tau\setminus\{r\}$.  At $x=0$
they are well defined by \Cref{lem:hard-feas}.
\end{theorem}

At $x=0$, the estimate is already contained in
\cite[Theorem~5.10 and Lemma~5.13]{CLMM2023}; it is also recovered below by
continuity.  To obtain the interval $[0,1]$, we extend the same tree argument
to positive activities.  For $0<x<1$, the soft recursion is a weighted-colour
recursion whose scaled cavity messages are subdistributions.  Four units of
palette slack below the root give a uniform contraction.  The contraction
implies tree total-influence decay (TID) and relative strong spatial mixing
(SSM) after a fixed depth; the root factors are handled by two exact
cancellations.
The large-girth transfer of~\cite{CLMM2023} then applies, and finite-state
continuity supplies the value at $x=0$.

\subsection{Weighted tree recursion}
\label{sec:hg-weighted-recursion}

We first take $0<x<1$ and put $s:=1-x$.  Root the free component of a tree
of maximum degree at most \(\Deg\) and condition on a fixed pinning, treating
every pinned neighbour as a boundary constraint.  At a free vertex $v$, let
$d_v$ be the number of free children and $k_v$ the number of pinned
neighbours.  As in \Cref{sec:potts-setup}, $\bcount_v^\tau(c)$ counts those of
colour $c$.  Thus
\[
  d_v+k_v\le\Deg-1
\]
away from the root, while $d_v+k_v\le\Deg$ at the root.  Define
\begin{equation}
  a_v(c)=x^{\bcount_v^\tau(c)},\qquad
  A_v=\sum_{c=1}^q a_v(c),\qquad
  \underline A_v=q-sk_v.
  \label{eq:hg-palette-mass}
\end{equation}
Bernoulli's inequality gives
\begin{equation}
  A_v\ge\underline A_v.
  \label{eq:hg-palette-lower}
\end{equation}
If $p_i$ is the cavity marginal of free child $i$, put $m_i=sp_i$.
The parent recursion is
\begin{equation}
  p_v(c)=
  \frac{a_v(c)\prod_{i=1}^{d_v}(1-m_i(c))}
       {\sum_{b=1}^q a_v(b)\prod_{i=1}^{d_v}(1-m_i(b))},
  \qquad m_v=sp_v.
  \label{eq:hg-scaled-recursion}
\end{equation}
Every $m_i$ is a subdistribution, meaning a nonnegative vector of total
mass at most one; here its total mass is exactly $s$.  More importantly,
for every non-root vertex,
\begin{equation}
  \underline A_v-d_v
    =q-sk_v-d_v
    \ge q-k_v-d_v\ge4.
  \label{eq:hg-four-slack}
\end{equation}
At the root the same calculation leaves three units of slack.  As in the hard
proof of~\cite{CLMM2023}, contraction is required only below the root; the
first step is retained as a bounded, possibly noncontracting map.

To integrate the Jacobian between two finite messages, we need a convex
domain that contains every genuine non-root message and every interpolation
point used in potential coordinates.  Put
\begin{equation}
 \mathcal D_s
 :=\left\{m\in[0,1/4]^q:
              \sum_{c\in\colours}m(c)\le s\right\}.
 \label{eq:hg-message-domain}
\end{equation}
This is a convex down-set.  The pointwise bound below shows that every
genuine non-root scaled message belongs to \(\mathcal D_s\) and has
total mass exactly \(s\); the extra points in
\eqref{eq:hg-message-domain} occur only along a potential-coordinate
interpolation.

We use the elementary product inequality
\begin{equation}
  0<B<1,\quad 0\le y_j\le B,\quad \sum_jy_j\le1
  \quad\Longrightarrow\quad
  \prod_j(1-y_j)\ge(1-B)^{1/B}.
  \label{eq:hg-product-lemma}
\end{equation}
Indeed, concavity of $\log(1-y)$ on the chord from $0$ to $B$ gives
$\log(1-y)\ge y\log(1-B)/B$.  Since the coefficient is negative,
\[
 \sum_j\log(1-y_j)
 \ge\frac{\log(1-B)}B\sum_jy_j
 \ge\frac{\log(1-B)}B,
\]
which proves the claim after exponentiation.

At a non-root vertex, inputs
\(m_i\in\mathcal D_s\) already give the crude invariant bound
\begin{equation}
\begin{aligned}
 Z_v
 &:=\sum_b a_v(b)\prod_i(1-m_i(b))\\
 &\ge A_v-\sum_i\sum_b a_v(b)m_i(b)
 \ge A_v-d_v\ge\underline A_v-d_v\ge4.
\end{aligned}
\label{eq:hg-Z-four}
\end{equation}
Consequently \(p_v(c)\le1/4\) and \(m_v=sp_v\) belongs to
\(\mathcal D_s\).  This proves invariance by induction from the leaves;
the sharper estimate below will also provide the coordinate bounds used in
the weighted norm.

\subsection{Uniform contraction below the root}
\label{sec:hg-nonroot-contraction}

We choose
\begin{equation}
  \phi(t)=2\operatorname{arctanh}\sqrt t,
  \qquad
  \phi'(t)=\frac1{\sqrt t(1-t)}
  \label{eq:hg-potential}
\end{equation}
because its derivative factorizes the transformed Jacobian into one factor
depending on the outgoing marginal and one measuring the incoming occupancy
of a single colour.  In the original probability coordinate the same
potential is \(\psi_s(p)=\phi(sp)\).  The map \(\phi\) is
increasing and concave on \([0,1/4]\); write
\[
  \boldsymbol\phi(m):=(\phi(m(c)))_{c\in\colours}.
\]
For the local map \((m_1,\ldots,m_{d_v})\mapsto p_v\), define
\[
 J^p_{v,i}
 :=D\boldsymbol\phi(p_v)\,
     \frac{\partial p_v}{\partial m_i}\,
     \bigl(D\boldsymbol\phi(m_i)\bigr)^{-1},
 \qquad
 J^p_v:=(J^p_{v,1}\mid\cdots\mid J^p_{v,d_v}).
\]
The superscript \(p\) indicates that this auxiliary Jacobian has the
unscaled output \(p_v\) and the scaled inputs \(m_i\).
Let \(\widehat J_{v,i}\) and
\(\widehat J_v=(\widehat J_{v,1}\mid\cdots\mid\widehat J_{v,d_v})\)
denote the analogous blocks with the scaled output \(\boldsymbol\phi(m_v)\).  At a zero
coordinate these matrices mean their continuous extensions from positive
coordinates.

For a block Jacobian \(J=(J_1,\ldots,J_d)\), a positive parent weight
\(w\), and positive child weights \(w_i\), put
\begin{equation}
 \norm{J}_{\mathbf w}
 =\sup_{(h_i)\ne0}
   \frac{w\norm{\sum_iJ_ih_i}_2}
        {\max_i w_i\norm{h_i}_2}.
 \label{eq:hg-weighted-norm}
\end{equation}
For \(d=0\) this empty block norm is zero.
Below, \(\norm{J^p_v}_{2\to2}\) denotes the spectral norm of the horizontally
concatenated \(q\)-by-\(d_vq\) block Jacobian.

\begin{lemma}[Exact Jacobian factorization]\label{lem:hg-jacobian-factor}
Write \(\Sigma_v(c)=\sum_i m_i(c)\) and
\[
  F_{1,v}=\max_{c\in\colours}\bigl(1-p_v(c)\bigr)^{-2},
  \qquad
  F_{2,v}=\max_{c\in\colours}p_v(c)\Sigma_v(c).
\]
Then \(\norm{J^p_v}_{2\to2}^2\le F_{1,v}F_{2,v}\).
\end{lemma}

\begin{proof}
Write \(p=p_v\) and \(\gamma(c)=\sqrt{p(c)}/(1-p(c))\).
Differentiating the \(\phi\)-transformed recursion gives
\[
  (J^p_{v,i})_{c,b}
   =\gamma(c)\bigl(p(b)-\one{b=c}\bigr)\sqrt{m_i(b)}.
\]
Hence, for a unit vector \(h\),
\[
  \norm{J^p_v}_{2\to2}^2
   =\max_{\norm{h}_2=1}\sum_c\Sigma_v(c)
      \bigl(\gamma(c)h(c)-p(c)\langle\gamma,h\rangle\bigr)^2 .
\]
Set \(u(c)=\gamma(c)h(c)\) and
\(U_1=\sum_cu(c)=\langle\gamma,h\rangle\).  The constraint
\(\norm h_2=1\) becomes
\[
  \sum_c\frac{u(c)^2(1-p(c))^2}{p(c)}=1.
\]
Since \(p(c)\Sigma_v(c)\le F_{2,v}\), the identity
\[
  \sum_c\frac{(u(c)-p(c)U_1)^2}{p(c)}
  =\sum_c\frac{u(c)^2}{p(c)}-U_1^2
\]
gives
\[
\begin{aligned}
 \sum_c\Sigma_v(c)(u(c)-p(c)U_1)^2
 &\le F_{2,v}\sum_c\frac{(u(c)-p(c)U_1)^2}{p(c)}\\
 &\le F_{2,v}\sum_c\frac{u(c)^2}{p(c)}
 \le F_{1,v}F_{2,v}.
\end{aligned}
\]
\end{proof}

We next bound the two factors.  The quantity \(B_v\) controls each
coordinate of \(p_v\), while \(\zeta_v\) records the loss in the occupancy
estimate.

\begin{lemma}[Pointwise marginal bound]
\label{lem:hg-continuous-marginal}
Consider \eqref{eq:hg-scaled-recursion} at a non-root vertex \(v\) with
$m_i\in\mathcal D_s$ for every free child.  With $d=d_v$ and
\(\underline A=\underline A_v\), define
\begin{equation}
  \xi_v=\left(\frac43\right)^{4d/(\underline A-1)},
  \qquad
  B_v=\frac{\xi_v}{\underline A-1+\xi_v}.
  \label{eq:hg-xi-B}
\end{equation}
Then
\begin{equation}
  \frac{p_v(c)}{1-p_v(c)}
    \le\frac{\xi_v}{\underline A-1},
  \qquad p_v(c)\le B_v\le\frac14,
  \qquad m_v(c)\le\frac14
  \label{eq:hg-marginal-bound}
\end{equation}
for every colour $c$.  In particular,
\(m_v\in\mathcal D_s\).
\end{lemma}

\begin{proof}
The crude estimate \eqref{eq:hg-Z-four} already gives
\(p_v(c)\le1/4\).  For the sharper recursive estimate, the claim is
immediate if \(a_v(c)=0\).  Otherwise set
\(S=A_v-a_v(c)\ge\underline A-1\).  Weighted AM--GM and
\(a_v(b)\le1\) give
\begin{align*}
 \sum_{b\ne c}a_v(b)\prod_i(1-m_i(b))
 &\ge S\prod_{b\ne c}
       \left(\prod_i(1-m_i(b))\right)^{a_v(b)/S}\\
 &\ge S\prod_i
       \left(\prod_{b\ne c}(1-m_i(b))\right)^{1/S}.
\end{align*}
By \eqref{eq:hg-product-lemma} with $B=1/4$, the inner product is at
least $(3/4)^4$.  The function
\(S\mapsto S(3/4)^{4d/S}\) is increasing, so
\[
  \frac{p_v(c)}{1-p_v(c)}
  \le \frac1{S(3/4)^{4d/S}}
  \le \frac{\xi_v}{\underline A-1}.
\]
Finally \(\underline A\ge d+4\) implies
\(\xi_v/(\underline A-1)\le1/3\).  For $d=0$ this is immediate; for $d\ge1$
it reduces to
$(4/3)^{4d/(d+3)}\le(d+3)/3$, whose logarithmic derivative is
nonnegative after equality at $d=1$.  This proves $B_v\le1/4$ and hence all
of \eqref{eq:hg-marginal-bound}.  The induction starts at a free leaf, where
the same argument has $d=0$.
\end{proof}

For every non-root vertex $y$, define its local $\xi_y,B_y$ by
\eqref{eq:hg-xi-B}, with $v$ replaced by $y$, and put
\begin{equation}
  \zeta_y=
  \left(\frac{\underline A_y-1}{\xi_y}+1\right)
  \log\left(1+\frac{\xi_y}{\underline A_y-1}\right)-1.
  \label{eq:hg-zeta}
\end{equation}
For a child \(y_i\) indexed by \(i\), abbreviate the corresponding
quantities by \(B_i:=B_{y_i}\) and \(\zeta_i:=\zeta_{y_i}\).

\begin{lemma}[Occupancy-factor bound]
\label{lem:hg-continuous-amortized}
Under the hypotheses of \Cref{lem:hg-continuous-marginal}, suppose in addition
that each free-child input satisfies $m_i(c)\le B_i$ for all $c$.  Then, for
every colour $c$,
\begin{equation}
  p_v(c)\sum_{i=1}^{d_v}m_i(c)
  \le \frac1{e\underline A_v}
       \exp\left\{
          \frac1{\underline A_v}\sum_{i=1}^{d_v}(\zeta_i+1)
       \right\}.
  \label{eq:hg-amortized}
\end{equation}
\end{lemma}

\begin{proof}
For fixed $c$, the $(d_v+1)$-term AM--GM inequality gives
\[
 a_v(c)\prod_i(1-m_i(c))\sum_i m_i(c)\le e^{-1}.
\]
The definition gives $B_i>0$, while the same slack calculation as in
\Cref{lem:hg-continuous-marginal} gives $B_i\le1/4$.
Writing
\(Z_v=\sum_b a_v(b)\prod_i(1-m_i(b))\)
for the recursion denominator, weighted AM--GM, \(a_v(b)\le1\), and
\eqref{eq:hg-product-lemma} give
\begin{align*}
 Z_v
 &\ge A_v\prod_b
    \left(\prod_i(1-m_i(b))\right)^{a_v(b)/A_v}\\
 &\ge A_v\prod_i
    \left(\prod_b(1-m_i(b))\right)^{1/A_v}\\
 &\ge A_v\prod_i(1-B_i)^{1/(A_vB_i)}\\
 &\ge \underline A_v
       \prod_i(1-B_i)^{1/(\underline A_vB_i)}.
\end{align*}
The last step uses that $t\mapsto t\exp(-H/t)$ is increasing for fixed
$H\ge0$.  Since
$B_i^{-1}\log(1/(1-B_i))=\zeta_i+1$, division proves
\eqref{eq:hg-amortized}.
\end{proof}

\begin{lemma}[Weighted Jacobian contraction]
\label{lem:hg-strong-jacobian}
At each non-root vertex set
$w_v=(1-\zeta_v)^{-1/2}$, and for a child \(y_i\) write
\(w_i:=w_{y_i}\).  For every $0<x<1$, every fixed pinning, and
every collection of free-child inputs $m_i\in\mathcal D_s$ satisfying
$m_i(c)\le B_i$, the transformed recursion in
\eqref{eq:hg-scaled-recursion} satisfies
\begin{equation}
  \norm{\widehat J_v}_{\mathbf w}\le
  \rho:=\exp\left(-\frac4{81q}\right)<1.
  \label{eq:hg-rho}
\end{equation}
Moreover
\begin{equation}
  1\le w_v\le e^{1/12},
  \label{eq:hg-weight-range}
\end{equation}
\end{lemma}

\begin{proof}
First reduce to the unscaled output.  Regard \(p_v\) in
\eqref{eq:hg-scaled-recursion} as a recursion with weighted colour
availability \(a_v(c)\) and subdistribution inputs \(m_i\).
The calculation below is first made in the relative interior of
$\prod_i\mathcal D_s$; the displayed formulas extend continuously to
the whole closed domain, including when an input coordinate tends to zero.
The chain rule with \eqref{eq:hg-potential} gives
\begin{equation}
 \widehat J_{v,i}
 =\operatorname{diag}\left(
   \frac{\sqrt s(1-p_v(c))}{1-sp_v(c)}
  \right)_{c=1}^q J^p_{v,i}.
 \label{eq:hg-scaled-row-factor}
\end{equation}
The diagonal factor is at most one because
\[
  1-st-\sqrt s(1-t)
   =(1-\sqrt s)(1+\sqrt s\,t)\ge0
  \qquad(0\le t\le1).
\]
Thus it is enough to bound the unscaled Jacobian with
\Cref{lem:hg-jacobian-factor}.  That factorization uses only
\(\sum_cp_v(c)=1\) and \(\Sigma_v(c)\ge0\), not \(\sum_bm_i(b)=1\); the weighted
colour availability enters only through \(Z_v\), controlled by
\Cref{lem:hg-continuous-marginal,lem:hg-continuous-amortized}.
Bounding \(F_{1,v}\le
\exp\bigl(2\xi_v/(\underline A_v-1)\bigr)\) by
\Cref{lem:hg-continuous-marginal} and \(F_{2,v}\) by
\eqref{eq:hg-amortized}, and using
\(d_v\le\underline A_v-4\) to produce the term
\(-4/\underline A_v\), yields
\begin{equation}
 \norm{J^p_v}_{2\to2}^2
 \le\frac1{\underline A_v}
 \exp\left\{
   \frac{2\xi_v}{\underline A_v-1}-\frac4{\underline A_v}
   +\frac1{\underline A_v}\sum_i\zeta_i
 \right\}.
 \label{eq:hg-euclidean-jacobian}
\end{equation}

Let \(y_v=\xi_v/(\underline A_v-1)\le1/3\).  The elementary entropy inequality
\cite[Fact~B.5]{CLMM2023} used in the weighted colouring recursion gives
\begin{equation}
 0<\zeta_v<1,\qquad
 \frac1{1-\zeta_v}\le e^{y_v/2},
 \label{eq:hg-entropy-inequality}
\end{equation}
and therefore \eqref{eq:hg-weight-range}.  Directly from
\eqref{eq:hg-weighted-norm},
\[
 \norm{J}_{\mathbf w}^2
 \le \norm{J}_{2\to2}^2 w_v^2\sum_iw_i^{-2}.
\]
Here \(w_i^{-2}=1-\zeta_i\), so
\(\sum_iw_i^{-2}=d_v-\sum_i\zeta_i\).
Applying this bound to \eqref{eq:hg-euclidean-jacobian} gives
\begin{align}
 \norm{\widehat J_v}_{\mathbf w}^2
 &\le
 \exp\left\{
   \frac{5\xi_v}{2(\underline A_v-1)}-\frac8{\underline A_v}
 \right\}.
 \label{eq:hg-weighted-jacobian-bound}
\end{align}
Indeed, with \(\bar\zeta_v=\underline A_v^{-1}\sum_i\zeta_i\),
\[
 \left(\frac{d_v}{\underline A_v}-\bar\zeta_v\right)e^{\bar\zeta_v}
 \le\frac{d_v}{\underline A_v}
 \le1-\frac4{\underline A_v}
 \le e^{-4/\underline A_v}.
\]
It remains to retain a strict margin.  Since
\(d_v\le\underline A_v-4\),
\[
 \xi_v\le\frac{256}{81}
 \exp\left(-\frac{12\log(4/3)}{\underline A_v-1}\right).
\]
Writing \(u=(\underline A_v-1)^{-1}\) and using
$\log(1+u)\le u\le12\log(4/3)u$ shows
\[
 \xi_v\frac{\underline A_v}{\underline A_v-1}\le\frac{256}{81}.
\]
Consequently the exponent in \eqref{eq:hg-weighted-jacobian-bound} is at
most \(-8/(81\underline A_v)\le-8/(81q)\).  Taking square roots proves
\eqref{eq:hg-rho}.  If $d_v=0$, the Jacobian is zero and the conclusion is
immediate.
\end{proof}

\begin{lemma}[From Jacobian contraction to message contraction]
\label{lem:hg-finite-difference}
Fix the local colour weights at a non-root vertex \(v\) and the child structural
parameters that determine \(B_i,w_i\).  Let
$(m_i)_{i=1}^{d_v}$ and $(m'_i)_{i=1}^{d_v}$ be two collections in
$\mathcal D_s$ such that
$m_i(c),m'_i(c)\le B_i$ for every $i,c$, and let
\(m_v,m'_v\) be their respective outputs under
\eqref{eq:hg-scaled-recursion}.  Then
\begin{equation}
 w_v\norm{\boldsymbol\phi(m_v)-\boldsymbol\phi(m'_v)}_2
 \le
 \rho\max_i
 w_i\norm{\boldsymbol\phi(m_i)-\boldsymbol\phi(m'_i)}_2 .
 \label{eq:hg-finite-difference}
\end{equation}
Both sides are interpreted as zero when $d_v=0$.
\end{lemma}

\begin{proof}
When $d_v=0$, the two outputs are identical.  Otherwise, for
$t\in[0,1]$ define the potential-coordinate segment
\[
 u_i(t)=(1-t)\boldsymbol\phi(m_i)+t\boldsymbol\phi(m'_i),
 \qquad
 m_i(t)=\boldsymbol\phi^{-1}(u_i(t)).
\]
Write \(u(t)=(u_i(t))_{i=1}^{d_v}\) and
\(m(t)=(m_i(t))_{i=1}^{d_v}\); thus
\(u'(t)=(\boldsymbol\phi(m'_i)-\boldsymbol\phi(m_i))_{i=1}^{d_v}\).
Since $\phi$ is increasing and concave on $[0,1/4]$, its inverse is
increasing and convex on $[0,\phi(1/4)]$.  Coordinatewise,
\[
 m_i(t,c)
 \le (1-t)m_i(c)+tm'_i(c)\le B_i,
 \qquad
 \sum_c m_i(t,c)
 \le (1-t)\sum_c m_i(c)+t\sum_c m'_i(c)
 \le s.
\]
Thus the entire segment lies in $\mathcal D_s$ and satisfies the sharper
coordinate bounds needed in \Cref{lem:hg-strong-jacobian}.  The fundamental
theorem of calculus and \eqref{eq:hg-rho} now give
\begin{align*}
 w_v\norm{\boldsymbol\phi(m_v)-\boldsymbol\phi(m'_v)}_2
 &\le
 \int_0^1
  w_v\norm{\widehat J_v(m(t))
                \bigl(u'(t)\bigr)}_2\,dt\\
 &\le
 \rho\max_iw_i\norm{\boldsymbol\phi(m_i)-\boldsymbol\phi(m'_i)}_2.
\end{align*}
If an endpoint has a zero coordinate, approximate it from the relative
interior and use continuity.
\end{proof}

\subsection{Total-influence decay and spatial mixing}
\label{sec:hg-tree-consequences}

We first fix the orientation and the messages used in the factorization.
Fix a tree \(\mathcal T=(V_{\mathcal T},E_{\mathcal T})\) of maximum
degree at most \(\Deg\), a set
\(\Lambda\subseteq V_{\mathcal T}\), a pinning
\(\tau\colon\Lambda\to\colours\), and a root
\(r\in V_{\mathcal T}\setminus\Lambda\).
Let \(\mathcal T_\tau\) be the component of \(r\) in
\(\mathcal T-\Lambda\), and let
\(\operatorname{dist}_{\mathcal T_\tau}\) denote its graph distance.
Root \(\mathcal T_\tau\) at \(r\).  For every
\(y\in V(\mathcal T_\tau)\setminus\{r\}\), let \(y^-\) be its parent and
let \(\mathcal T_{y\to y^-}\) be the component containing \(y\) after
the edge \(yy^-\) is removed from \(\mathcal T_\tau\).  Write
\[
  p_{y\to y^-}
\]
for the cavity marginal of $y$ in this component, including all fixed
pinnings incident to vertices of the component, and put
$m_{y\to y^-}=sp_{y\to y^-}$.  At the root, $p_r$ is instead the
ordinary marginal in the full rooted component and $m_r=sp_r$.  Thus
the root object is not being asserted to satisfy the non-root marginal bound.
By \Cref{lem:hg-continuous-marginal}, every cavity message
$m_{y\to y^-}$ does satisfy
\begin{equation}
  0\le m_{y\to y^-}(c)\le\frac14,
  \label{eq:hg-cavity-quarter}
\end{equation}
because deleting the parent edge leaves at most $\Deg-1$ child and pinned
constraints at $y$.

For a parent $y$ and its free child $w$, let $J_{y,w}$ denote the block
Jacobian, in the potential $\psi_s(p)=\phi(sp)$, of the appropriate
oriented update with respect to the incoming message from $w$: it outputs
the full marginal $p_r$ when $y=r$, and it outputs the cavity marginal
$p_{y\to y^-}$ when $y\ne r$.  In the scaled coordinate this is one block
of \(\widehat J_y\) when $y\ne r$, and is then controlled by
\Cref{lem:hg-strong-jacobian}; the root block is handled below.  Define,
with the same distinction between the root and non-root messages,
\begin{equation}
  D_r=\operatorname{diag}\bigl(p_r\psi_s'(p_r)\bigr)
     =\operatorname{diag}\frac{\sqrt{m_r}}{1-m_r},
  \qquad
  D_{y\to y^-}
     =\operatorname{diag}\frac{\sqrt{m_{y\to y^-}}}
                                  {1-m_{y\to y^-}} .
  \label{eq:hg-D-matrix}
\end{equation}
Products such as \(p_r\psi_s'(p_r)\) are coordinatewise.
Let \(I_q\) denote the \(q\)-by-\(q\) identity matrix and
\(\mathbf1_q\in\mathbb R^q\) the all-ones column vector.

\begin{lemma}[Bounded root response]\label{lem:hg-root-response}
For every free child \(u\) of the root,
\begin{equation}
 D_r^{-1}J_{r,u}
 =-\bigl(I_q-\mathbf1_qp_r^{\mathsf T}\bigr)
   \operatorname{diag}\bigl(\sqrt{m_{u\to r}}\bigr),
 \qquad
 \norm{D_r^{-1}J_{r,u}}_{2\to2}
 \le K_q:=\frac{1+\sqrt q}{2}.
 \label{eq:hg-first-jacobian}
\end{equation}
Thus the first influence factor is bounded without a lower bound on any
coordinate of \(p_r\).
\end{lemma}

\begin{proof}
The raw derivative of \eqref{eq:hg-scaled-recursion} is
\[
 \frac{\partial p_r(c)}{\partial p_{u\to r}(b)}
 =\frac{s p_r(c)
        [p_r(b)-\one{b=c}]}{1-m_{u\to r}(b)},
 \qquad
 \psi_s'(p_{u\to r}(b))^{-1}
 =\frac{\sqrt{m_{u\to r}(b)}
              (1-m_{u\to r}(b))}{s}.
\]
The factor \(p_r(c)\) cancels against \(D_r^{-1}\) before any norm is
taken, giving the displayed identity.  By
\eqref{eq:hg-cavity-quarter},
\(\norm{\operatorname{diag}\sqrt{m_{u\to r}}}_{2\to2}\le1/2\), while
\(\norm{I_q-\mathbf1_qp_r^{\mathsf T}}_{2\to2}\le1+\sqrt q\).
\end{proof}

For this pinning \(\tau\colon\Lambda\to\colours\), write
\(\mu_v^\tau\) for the one-site marginal
at a free vertex \(v\) under the corresponding tree Gibbs law and
\(\mu_v^{\tau^{r=a}}\) for that marginal after additionally pinning a free
vertex \(r\) to \(a\).  For distinct free vertices \(r,v\), let
\(\Psi_{r,v}\) be the \(q\)-by-\(q\) influence block
\[
 \Psi_{r,v}(a,c)
 :=\mu_v^{\tau^{r=a}}(c)-\mu_v^\tau(c);
\]
we suppress its dependence on \(\tau\).  If deleting the vertices pinned by
\(\tau\) separates \(r\) from \(v\), this block is zero.  Define the
level-\(\ell\) total influence by
\begin{equation}
 \mathcal I_\ell(r;\tau)
 :=\max_{a,b\in\colours}
   \sum_{\substack{v\in V(\mathcal T_\tau)\\
                   \operatorname{dist}_{\mathcal T_\tau}(r,v)=\ell}}
   \norm{\mu_v^{\tau^{r=a}}-\mu_v^{\tau^{r=b}}}_{\rm TV}.
 \label{eq:hg-tid-definition}
\end{equation}
If \(r\) and \(v\) remain connected, let their free path be
\[
  r=u_0,u_1,\ldots,u_\ell=v.
\]
The influence/Jacobian identity of \cite[Lemma~8.7]{CLMM2023}, applied in
the successive rooted subtrees, gives
\begin{equation}
  \Psi_{r,v}
   =D_r^{-1}J_{u_0,u_1}
      \Bigl(\prod_{j=1}^{\ell-1}J_{u_j,u_{j+1}}\Bigr)
      D_{v\to u_{\ell-1}}
      \bigl(I_q-\mathbf1_qp_{v\to u_{\ell-1}}^{\mathsf T}\bigr).
  \label{eq:hg-psi-factorization}
\end{equation}
The product is ordered away from the root, is empty when \(\ell=1\), and
the intermediate diagonal matrices telescope.  The terminal factors use
the \emph{cavity} law in
\(\mathcal T_{v\to u_{\ell-1}}\).  The identity is exact, and
\Cref{lem:hg-root-response} shows that its first factor extends boundedly
to a zero root coordinate.

\begin{lemma}[Tree total-influence decay]
\label{lem:hg-tid}
Uniformly over $0<x<1$, every tree
\(\mathcal T=(V_{\mathcal T},E_{\mathcal T})\) of maximum degree at most
\(\Deg\), every pinning \(\tau\colon\Lambda\to\colours\) with
\(\Lambda\subseteq V_{\mathcal T}\), every root
\(r\in V_{\mathcal T}\setminus\Lambda\), and every integer $\ell\ge1$ satisfy
\begin{equation}
 \mathcal I_\ell(r;\tau)
 \le C_0\rho^{\ell-1}
 =C_{\rm TID}\rho^\ell,
\label{eq:hg-tid}
\end{equation}
where one may take
\begin{equation}
 C_0=\frac{\Deg e^{1/12}\sqrt q(1+\sqrt q)^2}{3},
 \qquad
 C_{\rm TID}=\frac{C_0}{\rho}.
 \label{eq:hg-tid-constant}
\end{equation}
The level sum ranges only over free vertices.
No positive lower bound on a one-site marginal enters this constant.
\end{lemma}

\begin{proof}
The first factor is bounded by \Cref{lem:hg-root-response}; it need not
contract.  At the terminal end, \eqref{eq:hg-cavity-quarter} applies to
$m_{v\to u_{\ell-1}}$, and therefore
\begin{equation}
 \norm{D_{v\to u_{\ell-1}}
       (I_q-\mathbf1_qp_{v\to u_{\ell-1}}^{\mathsf T})}_{2\to2}
 \le\frac{2(1+\sqrt q)}3,
 \label{eq:hg-terminal-bound}
\end{equation}
because $\norm{D_{v\to u_{\ell-1}}}_{2\to2}\le2/3$ and
$\norm{I_q-\mathbf1_qp_{v\to u_{\ell-1}}^{\mathsf T}}_{2\to2}\le1+\sqrt q$.

We now use \eqref{eq:hg-psi-factorization} simultaneously over a level.
If the level contains no free vertex, the claimed influence bound is
immediate.  Otherwise, group the terminals by their first child of $r$,
keep the corresponding
first factor \eqref{eq:hg-first-jacobian} intact, and peel the remaining
$\ell-1$ levels with the block norm of
\Cref{lem:hg-strong-jacobian}.  Every update peeled in this way is a
non-root cavity update.  Iterating the weighted inequality gives
$\rho^{\ell-1}$; the ratio between its terminal and initial weights is at
most $e^{1/12}$ by \eqref{eq:hg-weight-range}.  Finally sum the at most
$\Deg$ first-child blocks by the triangle inequality.  Thus, for arbitrary
terminal block vectors $h_v\in\mathbb R^q$,
\[
 \frac{
  \norm{\sum_{\operatorname{dist}_{\mathcal T_\tau}(r,v)=\ell}
              \Psi_{r,v}h_v}_2}
  {\max_v\norm{h_v}_2}
 \le
 \Deg e^{1/12}K_q\frac{2(1+\sqrt q)}3\rho^{\ell-1}.
\]
Here the product and the contraction exponent are interpreted as empty and
$1$, respectively, when $\ell=1$.  The standard block
$L^2$-to-$L^\infty$ duality conversion costs at most $\sqrt q$.  Indeed, if
\[
 \mathcal A(h):=
 \sum_{\operatorname{dist}_{\mathcal T_\tau}(r,v)=\ell}
       \Psi_{r,v}h_v ,
\]
the triangle
inequality through the unconditioned law makes the two factors \(1/2\) in
total variation cancel, and hence
\[
 \mathcal I_\ell(r;\tau)
 \le \max_a
       \sum_{\operatorname{dist}_{\mathcal T_\tau}(r,v)=\ell}\sum_c
       \abs{\Psi_{r,v}(a,c)}
 \le
 \sup_{\max_v\norm{h_v}_\infty\le1}\norm{\mathcal A(h)}_\infty
 \le
 \sqrt q\,
 \sup_{\max_v\norm{h_v}_2\le1}\norm{\mathcal A(h)}_2 .
\]
Using the definition of $K_q$ and writing
$\rho^{\ell-1}=\rho^{-1}\rho^\ell$ gives exactly
\eqref{eq:hg-tid}--\eqref{eq:hg-tid-constant}.
\end{proof}

Uniform relative SSM cannot start at distance one. Indeed, on a single
edge, compare boundary colours \(a\) and \(b\ne a\).  The probability that
the free endpoint has colour \(a\) is respectively
\[
  \frac{x}{q-1+x}
  \qquad\text{and}\qquad
  \frac1{q-1+x},
\]
so the ratio is \(x\) (or \(1/x\) in the reverse comparison).  The
large-girth transfer only needs a relative estimate beyond a fixed depth.

\begin{lemma}[Relative SSM beyond a fixed depth]
\label{lem:hg-eventual-relative-ssm}
For
\begin{equation}
 C_{\rm rel}:=\frac58\Deg q\,e^{1/12}\log3,
 \label{eq:hg-relative-constant}
\end{equation}
the following holds for every $0<x<1$ on every tree
\(\mathcal T=(V_{\mathcal T},E_{\mathcal T})\) of maximum degree at
most $\Deg$.  Fix \(\Lambda\subseteq V_{\mathcal T}\), a pinning
\(\tau\colon\Lambda\to\colours\), and a root
\(v\in V_{\mathcal T}\setminus\Lambda\).  For an integer \(K\ge2\), let
\[
  S_K(v):=
  \{u\in V_{\mathcal T}:
       \operatorname{dist}_{\mathcal T}(u,v)=K\},
\]
and let
\[
 \omega,\omega'\colon S_K(v)\setminus\Lambda\longrightarrow
 \colours
\]
be two assignments of all unpinned shell vertices.  Thus
\(\tau\cup\omega\) and \(\tau\cup\omega'\) are the unions of partial maps with
disjoint domains.  The two systems have the same fixed pinning \(\tau\) and
differ only on the shell.  Write
\[
 p_v:=\mu_v^{\tau\cup\omega},
 \qquad
 p'_v:=\mu_v^{\tau\cup\omega'}.
\]
Here each right-hand side denotes the one-site marginal at \(v\) under the
indicated pinning.
Then, for every root colour $c$,
\begin{equation}
 \left|\log\frac{p_v(c)}{p'_v(c)}\right|
 \le C_{\rm rel}\rho^{K-2}.
 \label{eq:hg-eventual-relative-ssm}
\end{equation}
Define
\[
 K_1:=\min\{K\in\mathbb Z_{\ge2}:C_{\rm rel}\rho^{K-2}\le1\}.
\]
For every integer \(K\ge K_1\), the ratio-form relative-SSM error is at most
\(2C_{\rm rel}\rho^{K-2}\).
The case \(x=0\) will be added after the final coupling bound by
finite-state continuity; at \(x=1\) the left side is zero.
\end{lemma}

\begin{proof}
The pinned children adjacent to the root belong to the fixed pinning $\tau$,
so their colour weights \(a_v(c)\) agree.  Write
\[
 P_c=\prod_i(1-m_i(c)),\qquad
 p_v(c)=\frac{a_v(c)P_c}{Z},\qquad
 Z=\sum_b a_v(b)P_b,
\]
and use primes for the second condition.  The potentially tiny root colour
weight cancels exactly:
\begin{equation}
 \log\frac{p_v(c)}{p'_v(c)}
 =\log\frac{P_c}{P'_c}+\log\frac{Z'}Z.
 \label{eq:hg-palette-cancellation}
\end{equation}
Because $K\ge2$, every message entering the root is the output of a non-root
recursion and hence is bounded by $1/4$.  Therefore
\[
 \left|\log\frac{P_c}{P'_c}\right|
 \le\frac43\sum_i|m_i(c)-m'_i(c)|.
\]
Moreover telescoping the products gives
$|Z-Z'|\le\sum_i\norm{m_i-m'_i}_1$, while the root version of
\eqref{eq:hg-Z-four} gives $Z,Z'\ge3$.  Consequently
\begin{equation}
 \left|\log\frac{p_v(c)}{p'_v(c)}\right|
 \le\frac53\sum_i\norm{m_i-m'_i}_1.
 \label{eq:hg-relative-to-messages}
\end{equation}

The update adjacent to the shell is not compared by a Jacobian, because its
local colour weights may change with the shell assignment.  Its two output
messages nevertheless obey \Cref{lem:hg-continuous-marginal}.  For
\(1\le j\le K-1\), let \(\mathcal F_j\) be the free vertices at distance
\(j\) from \(v\) in the same component after the vertices pinned by
\(\tau\) are deleted.  Orient this component toward \(v\), and write
\(y^-\) for the parent of \(y\).  Let
\[
 E_j:=
 \max_{y\in\mathcal F_j}
 w_y\norm{\boldsymbol\phi(m_{y\to y^-})
             -\boldsymbol\phi(m'_{y\to y^-})}_2 .
\]
We set the maximum over an empty set to zero.
The potential diameter and \eqref{eq:hg-weight-range} give
\[
 E_{K-1}\le e^{1/12}\sqrt q\,\phi(1/4)
 =e^{1/12}\sqrt q\log3.
\]
The two boundary assignments pin the same set of shell vertices, so the two
recursions have the same \(d_y,k_y,B_y,w_y\).  Every update strictly inside
that shell has the same local colour weights on both sides.  Thus
\Cref{lem:hg-finite-difference} gives
\[
 E_j\le\rho E_{j+1}
 \qquad(1\le j\le K-2).
\]
These are exactly \(K-2\) contracting levels.  Since \(w_i\ge1\) and
\(\inf_{[0,1/4]}\phi'=8/3\), Cauchy--Schwarz gives, for each message entering
the root,
\[
 \norm{m_i-m'_i}_1
 \le\frac38\sqrt q\,
       \norm{\boldsymbol\phi(m_i)-\boldsymbol\phi(m'_i)}_2
 \le\frac38q\,e^{1/12}\log3\,\rho^{K-2}.
\]
Summing over at most \(\Deg\) root children in
\eqref{eq:hg-relative-to-messages} proves
\eqref{eq:hg-eventual-relative-ssm} with
\eqref{eq:hg-relative-constant}.

To convert this logarithmic estimate to the relative-error form of
\cite[Definition~5.7]{CLMM2023}, put
\(\varepsilon_K:=C_{\rm rel}\rho^{K-2}\).  By the definition of \(K_1\), for
\(K\ge K_1\),
\[
 \left|\frac{p_v(c)}{p'_v(c)}-1\right|
 \le e^{\varepsilon_K}-1\le2\varepsilon_K .
\]
Thus the usual relative-SSM estimate holds uniformly at every
$K\ge K_1$; no assertion is made at distance one.
\end{proof}

\subsection{Transfer to large-girth graphs}

The transfer of~\cite{CLMM2023} still applies when relative SSM begins at a
fixed depth \(K_0\).

\begin{lemma}[Large-girth transfer from tree decay]
\label{lem:hg-eventual-transfer}
Let \(q\ge1\) and \(\Deg\ge3\) be integers, and let \(J\subseteq(0,1]\).
Suppose there are
$0<C_{\mathsf{INFL}},C_{\mathsf{SM}}<\infty$, $\delta\in(0,1)$, and a fixed
positive integer $K_0$ such that, for every \(x\in J\) and uniformly over all
trees \(T\in\Gdeg\), all pinnings \(\tau\) of \(T\), and all free roots, the
Potts law \(\mupin{T}{\tau}{x}\) has level-\(\ell\) tree total influence at
most $C_{\mathsf{INFL}}(1-\delta)^\ell$ for every integer $\ell\ge1$, and
relative SSM in the ratio sense of \cite[Definition~5.7]{CLMM2023} at
distance $K$ at most $C_{\mathsf{SM}}(1-\delta)^K$ for every integer
$K\ge K_0$.  Then there are constants \(g_{\mathsf{tr}}\in\mathbb Z_{\ge3}\)
and \(C_{\mathsf{tr}}\ge0\), depending only on
\(q,\Deg,C_{\mathsf{INFL}},C_{\mathsf{SM}},\delta,K_0\), and in particular
not on \(x\) or \(J\), with the following property.  For every \(x\in J\),
every \(G\in\Gdeg\), every pinning \(\tau\) with
\(\operatorname{girth}(G^\tau)\ge g_{\mathsf{tr}}\), every free source vertex
\(u\in V^\tau\), and all \(b,c\in\colours\),
\begin{equation}
 \WHam\!\left(
   \mupin{G}{\tau^{u=b}}{x},
   \mupin{G}{\tau^{u=c}}{x}
 \right)\le C_{\mathsf{tr}}.
 \label{eq:hg-transfer-W1}
\end{equation}
Both laws live on the common free set \(V^\tau\setminus\{u\}\).
\end{lemma}

\begin{proof}
Fix \(x\in J\), and regard each pinned pair \((G,\tau)\) with \(G\in\Gdeg\) as
a pairwise spin system on its free graph \(G^\tau\), with spin set
\(\colours\), the Potts edge factors, and external fields
\[
  \lambda_{v,c}=x^{\bcount_v^\tau(c)}.
\]
The family of these systems is closed under further pinnings and deletion of
free vertices.  Pinning a free vertex $r$ to $c$ removes one free incidence at
each neighbour and adds $c$ to its boundary count.  Hence, for every remaining
vertex $v$,
\[
  \deg_{G^{\tau^{r=c}}}(v)
    +\sum_d\bcount_v^{\tau^{r=c}}(d)
  =\deg_{G^\tau}(v)+\sum_d\bcount_v^\tau(d)
  \le\Deg,
\]
whereas taking a subgraph only lowers this budget.  Since \(x>0\), every
configuration has positive weight, so every spin is globally feasible after
every pinning and the feasibility requirements of the pairwise transfer hold
trivially.  Conditioning the source \(u\) fixes the scalar
\(\lambda_{u,c}\), which disappears on normalization, and turns each incident
edge into the boundary factor in \(\tau^{u=c}\).  Thus the conditional
exterior laws supplied by the transfer are exactly the two laws in
\eqref{eq:hg-transfer-W1}.

The pairwise extension follows from \cite[Remark~5.11]{CLMM2023}, and the
relative input is the ratio form in
\cite[Definition~5.7]{CLMM2023}.  It remains to check that the proof tolerates
the restriction \(K\ge K_0\).  The radius, the cutting depth, and the girth
bound chosen below depend only on
\(q,\Deg,C_{\mathsf{INFL}},C_{\mathsf{SM}},\delta,K_0\), so they serve every
\(x\in J\) at once.

The cutting depth \(K\) is chosen after the influence radius \(R\).  The error
in the proof of \cite[Lemma~5.14]{CLMM2023} is
\begin{equation}
  2C_{\mathsf{SM}}(1-\delta)^K\Deg^R
  +C_{\mathsf{INFL}}(1-\delta)^R.
  \label{eq:hg-transfer-error}
\end{equation}
First choose $R$, subject to all lower bounds on $R$ in the cited proof, so
large that
\[
 C_{\mathsf{INFL}}(1-\delta)^R
 \le\frac1{16R\log\Deg}.
\]
After this $R$ is fixed, choose $K$ subject to \emph{all} lower bounds on the
cutting depth in \cite[Lemmas~5.14 and~5.20]{CLMM2023}, and also
\[
 K\ge K_0,\qquad K>R,\qquad
 K\ge\left\lceil\frac{\log C_{\mathsf{SM}}}{\delta}\right\rceil,\qquad
 2C_{\mathsf{SM}}(1-\delta)^K\Deg^R
 \le\frac1{16R\log\Deg}.
\]
Such a $K$ exists; with this choice \eqref{eq:hg-transfer-error} is at most
$1/(8R\log\Deg)$, which is the intermediate influence-decay threshold in
\cite[Lemma~5.14]{CLMM2023}.
Finally take the girth bound to be the maximum of every girth lower bound in
the cited lemmas and $2K+2$.  Relative SSM is invoked only at this chosen
cutting depth, so the restriction \(K\ge K_0\) causes no change later in the
argument.  Lemmas~5.14 and~5.20 of~\cite{CLMM2023} therefore supply the
intermediate decay condition, and their Lemma~5.13 gives
\eqref{eq:hg-transfer-W1} with
\[
 C_{\mathsf{tr}}=2\Deg^R.
\]
\end{proof}

\begin{proof}[Proof of \Cref{thm:high-girth-soft-ci}]
\Cref{lem:hg-tid,lem:hg-eventual-relative-ssm} give, for every \(0<x<1\),
tree TID and relative SSM beyond a fixed depth, with constants and rate
uniform in \(x\).  Apply \Cref{lem:hg-eventual-transfer} with
\[
 J=(0,1),\qquad
 C_{\mathsf{INFL}}=C_{\rm TID},\qquad
 C_{\mathsf{SM}}=2C_{\rm rel}\rho^{-2},\qquad
 \delta=1-\rho,\qquad K_0=K_1.
\]
It supplies \(g_{\mathsf{tr}}\) and \(C_{\mathsf{tr}}\), depending only on
\((\Deg,q)\), such that, for every \(0<x<1\), whenever $G^\tau$ has girth at
least $g_{\mathsf{tr}}$,
\begin{equation}
  \WHam\!\left(
    \mupin{G}{\tau^{r=a}}{x},
    \mupin{G}{\tau^{r=b}}{x}
  \right)\le C_{\mathsf{tr}}.
  \label{eq:hg-potts-transfer-W1}
\end{equation}

They live on the common finite state space
$\colours^{V^\tau\setminus\{r\}}$.  At \(x=0\), the hypothesis
\(q\ge\Deg+3\) implies hard feasibility for every normalized child by
\Cref{lem:hard-feas}, so both limiting normalizing constants are positive.
Their weights are polynomials in \(x\), and
\Cref{lem:finite-endpoint-closure} passes the same bound to \(x=0\).  At
\(x=1\) all remaining spins are independent and the distance is zero.
Taking \(g_*=\max\{3,g_{\mathsf{tr}}\}\) and
\(C_*=C_{\mathsf{tr}}\) proves the theorem.
\end{proof}
